\section{The main theorems without hypotheses}
\label{sec:main-fully-unconditional}

The main theorems of \cite{FHX} were proved above under two hypotheses at the
countable algebraically closed subfields $k_0$ of $k$: Lemma~5.5 of \cite{FHX}
(Lemma~\ref{lem:wu-local-positivity}) and, in characteristic zero, the envelope
property of smooth projectively normal subvarieties of projective space
(Theorem~\ref{thm:import-sg-ep-char0}). Both hold over every algebraically
closed field: the first by Theorem~\ref{thm:lpn-local-positivity-instance}, the
second by Theorem~\ref{thm:kn-smooth-envelope-char-zero}. For general models
the theorems also assume fibre constancy at those subfields, which holds over
every algebraically closed field (Theorem~\ref{thm:gl-fibre-constancy}). This
section states the resulting theorems; none of them has a hypothesis beyond
those of \cite{FHX}.

\begin{theorem}[Fang--Han--Xia, Theorem~A]
\label{thm:fu-variety-main}
\uses{thm:var-main,thm:lpn-local-positivity-instance,thm:kn-smooth-envelope-char-zero}
\lean{EnvelopeProperty.Variety.hasEnvelopeProperty_of_isUnibranch,EnvelopeProperty.Variety.isPSH_uscReg_iSup_of_isUnibranch}
\leanok
Let $k$ be an algebraically closed field, $X$ an integral unibranch projective
variety over $k$ and $\theta\in N^1(X)_\R$. Then $\mathrm{EP}_X(\theta)$ holds:
every nonempty family $(\varphi_i)\subseteq\PSH(X,\theta)$ that is bounded above
satisfies $(\sup_i\varphi_i)^\star\in\PSH(X,\theta)$.
\end{theorem}
\begin{proof}
Every countable algebraically closed subfield $k_0$ of $k$ is algebraically
closed, so Lemma~5.5 of \cite{FHX} holds over $k_0$
(Theorem~\ref{thm:lpn-local-positivity-instance}), and if $k_0$ has
characteristic zero, so does the envelope property of smooth projectively normal
subvarieties of $\mathbf P^M_{k_0}$ in the classes $\theta\,c_1(\mathcal O(1))$,
$\theta\geq0$ (Theorem~\ref{thm:kn-smooth-envelope-char-zero}). These are the
hypotheses of Theorem~\ref{thm:var-main}, which gives the conclusion.
\end{proof}

\begin{theorem}[Fang--Han--Xia, Theorem~A, with the stalk-wise definition of unibranch]
\label{thm:fu-variety-main-stalks}
\uses{thm:fu-variety-main,lem:var-unibranch-stalks}
\lean{EnvelopeProperty.Variety.hasEnvelopeProperty_of_forall_isUnibranchLocalRing_stalk,EnvelopeProperty.Variety.isPSH_uscReg_iSup_of_forall_isUnibranchLocalRing_stalk}
\leanok
Let $k$ be an algebraically closed field, $X$ an integral projective variety over
$k$ such that every local ring $\mathcal O_{X,x}$ is unibranch, and
$\theta\in N^1(X)_\R$. Then every nonempty family
$(\varphi_i)\subseteq\PSH(X,\theta)$ that is bounded above satisfies
$(\sup_i\varphi_i)^\star\in\PSH(X,\theta)$; that is, $\mathrm{EP}_X(\theta)$
holds.
\end{theorem}
\begin{proof}
By Lemma~\ref{lem:var-unibranch-stalks}, $X$ is unibranch in the sense of
Definition~\ref{def:var-unibranch}; apply Theorem~\ref{thm:fu-variety-main}.
\end{proof}

This is \texttt{thm:intro-main} of \cite{FHX} word for word, since
\cite[\S1.1]{FHX} defines an integral variety to be unibranch when every local
ring has one branch.

\begin{corollary}[Fang--Han--Xia, Corollary~1.1]
\label{cor:fu-variety-characterization}
\uses{cor:var-characterization,thm:lpn-local-positivity-instance,thm:kn-smooth-envelope-char-zero}
\lean{EnvelopeProperty.Variety.isUnibranch_tfae_hasEnvelopeProperty}
\leanok
Let $k$ be an algebraically closed field and $X$ an integral projective variety
over $k$. The following are equivalent:
\begin{enumerate}
\item $X$ is unibranch;
\item[(1$'$)] every local ring $\mathcal O_{X,x}$ is unibranch;
\item[(1$''$)] every local ring $\mathcal O_{X,x}$ is geometrically unibranch;
\item[(2)] $\mathrm{EP}_X(\theta)$ holds for every $\theta\in N^1(X)_\R$;
\item[(3)] $\mathrm{EP}_X(\omega)$ holds for some ample class $\omega\in\operatorname{Amp}(X)$.
\end{enumerate}
\end{corollary}
\begin{proof}
As in the proof of Theorem~\ref{thm:fu-variety-main}, the hypotheses of
Corollary~\ref{cor:var-characterization} hold at every countable algebraically
closed subfield of $k$, by Theorems~\ref{thm:lpn-local-positivity-instance}
and~\ref{thm:kn-smooth-envelope-char-zero}; that corollary gives the
equivalences.
\end{proof}

\begin{theorem}[Fang--Han--Xia, Theorem~A, for multiprojective subvarieties]
\label{thm:fu-quot-main}
\uses{thm:lpn-main-unconditional,thm:kn-smooth-envelope-char-zero}
\lean{EnvelopeProperty.MultiModel.hasEnvelopeProperty_quotModel_of_isUnibranch,EnvelopeProperty.MultiModel.isPSH_uscReg_iSup_quotModel_of_isUnibranch}
\leanok
Let $k$ be an algebraically closed field, $I\subseteq k[x^{(c)}_j]$ a prime
ideal relevant at $j_0$, and suppose that the quotient model $\mathcal X_I$ of
$X=V(I)\subseteq\mathbf P^{M_1}\times\dots\times\mathbf P^{M_r}$ is unibranch.
Then $\mathrm{EP}_{\mathcal X_I}(\theta)$ holds for every $\theta\in\R^r$:
every nonempty family $(\varphi_i)\subseteq\PSH(\mathcal X_I,\theta)$ bounded
above satisfies $(\sup_i\varphi_i)^\star\in\PSH(\mathcal X_I,\theta)$.
\end{theorem}
\begin{proof}
Theorem~\ref{thm:lpn-main-unconditional} gives the statement under the
envelope property of smooth projectively normal subvarieties of
$\mathbf P^M_{k_0}$ in characteristic zero at the countable algebraically closed
subfields $k_0$ of $k$; it holds by Theorem~\ref{thm:kn-smooth-envelope-char-zero}.
\end{proof}

\begin{corollary}[Fang--Han--Xia, Corollary~1.1, for multiprojective subvarieties]
\label{cor:fu-quot-characterization}
\uses{thm:lpn-main-unconditional,thm:kn-smooth-envelope-char-zero}
\lean{EnvelopeProperty.MultiModel.quotModel_isUnibranch_tfae}
\leanok
In the setting of Theorem~\ref{thm:fu-quot-main} without the unibranch
assumption, the following are equivalent: (1) $\mathcal X_I$ is unibranch;
(1$'$) every local ring of every affine chart of $\mathcal X_I$ is unibranch;
(1$''$) every such local ring is geometrically unibranch; (2)
$\mathrm{EP}_{\mathcal X_I}(\theta)$ for every $\theta\in\R^r$; (3)
$\mathrm{EP}_{\mathcal X_I}(\omega)$ for some $\omega\in\R^r$ with all
$\omega_c>0$.
\end{corollary}
\begin{proof}
Theorem~\ref{thm:lpn-main-unconditional} gives the equivalences under the
characteristic-zero envelope hypothesis at the countable algebraically closed
subfields of $k$, which holds by Theorem~\ref{thm:kn-smooth-envelope-char-zero}.
\end{proof}

\begin{theorem}[The main theorem for models {\cite[Theorem~A]{FHX}}]
\label{thm:fu-model-main}
\uses{thm:lpn-main-unconditional,thm:kn-smooth-envelope-char-zero,thm:gl-fibre-constancy,prop:gl-ample-countable}
\lean{EnvelopeProperty.MultiModel.hasEnvelopeProperty_of_isUnibranch,EnvelopeProperty.MultiModel.isPSH_uscReg_iSup_of_isUnibranch,EnvelopeProperty.MultiModel.hasEnvelopeProperty_of_isNormal,EnvelopeProperty.ProjSubvariety.hasEnvelopeProperty_one_of_isProjNormal_of_subfield}
\leanok
Let $k$ be an algebraically closed field. Then:
\begin{enumerate}
\item every projectively normal $X_0=V(I_0)\subseteq\mathbf P^M_{k_0}$ over a
countable algebraically closed subfield $k_0$ of $k$ has
$\mathrm{EP}_{X_0}(1)$;
\item every normal model $\mathcal X$ of rank $r$ over $k$ has
$\mathrm{EP}_{\mathcal X}(\theta)$ for every $\theta\in\R^r$;
\item every unibranch model $\mathcal X$ of rank $r$ over $k$ has
$\mathrm{EP}_{\mathcal X}(\theta)$ for every $\theta\in\R^r$: every nonempty
family $(\varphi_i)\subseteq\PSH(\mathcal X,\theta)$ bounded above satisfies
$(\sup_i\varphi_i)^\star\in\PSH(\mathcal X,\theta)$.
\end{enumerate}
\end{theorem}
\begin{proof}
Theorem~\ref{thm:lpn-main-unconditional} gives (1)--(3) under fibre constancy
and the characteristic-zero envelope hypothesis at the countable algebraically
closed subfields of $k$; (1) is Proposition~\ref{prop:gl-ample-countable}. Such
subfields are algebraically closed, so fibre constancy holds by
Theorem~\ref{thm:gl-fibre-constancy} and the envelope hypothesis by
Theorem~\ref{thm:kn-smooth-envelope-char-zero}.
\end{proof}

\begin{corollary}[Fang--Han--Xia, Corollary~8.4, for models]
\label{cor:fu-model-psh2-unibranch}
\uses{cor:nn-psh2-unibranch,thm:lpn-local-positivity-instance,thm:kn-smooth-envelope-char-zero,thm:gl-fibre-constancy}
\lean{EnvelopeProperty.MultiModel.isPSH_of_antitone_of_isUnibranch}
\leanok
Let $k$ be an algebraically closed field and $\mathcal X$ a unibranch model of
rank $r$ over $k$. If $\theta_i\to\theta$ in $\R^r$ and a net
$u_i\in\PSH(\mathcal X,\theta_i)$ decreases pointwise to
$u\not\equiv-\infty$, then $u\in\PSH(\mathcal X,\theta)$.
\end{corollary}
\begin{proof}
The hypotheses of Corollary~\ref{cor:nn-psh2-unibranch} at the countable
algebraically closed subfields of $k$ hold: Lemma~5.5 of \cite{FHX} by
Theorem~\ref{thm:lpn-local-positivity-instance}, fibre constancy by
Theorem~\ref{thm:gl-fibre-constancy}, and the characteristic-zero envelope
property by Theorem~\ref{thm:kn-smooth-envelope-char-zero}.
\end{proof}

\cite[Corollary~8.4]{FHX} is stated for an integral unibranch projective variety
$X$ and $\theta\in N^1(X)_\R$; Corollary~\ref{cor:fu-model-psh2-unibranch} is its
form on models. The variety form would follow by transporting decreasing nets and
the convergence of classes through a presentation
(Theorem~\ref{thm:var-psh-comparison}); it is not formalized.

\begin{corollary}[Characterization of unibranch models {\cite[Corollary~1.1]{FHX}}]
\label{cor:fu-model-characterization}
\uses{cor:nn-characterization,thm:lpn-local-positivity-instance,thm:kn-smooth-envelope-char-zero,thm:gl-fibre-constancy}
\lean{EnvelopeProperty.MultiModel.isUnibranch_tfae_hasEnvelopeProperty}
\leanok
Let $k$ be an algebraically closed field. For a model $\mathcal X$ of
rank $r$ over $k$ the following are equivalent: (1) $\mathcal X$ is unibranch;
(1$'$) all local rings of all chart algebras are unibranch; (1$''$) they are
all geometrically unibranch; (2) $\mathrm{EP}_{\mathcal X}(\theta)$ for every
$\theta\in\R^r$; (3) $\mathrm{EP}_{\mathcal X}(\omega)$ for some $\omega\in\R^r$
with all $\omega_c>0$.
\end{corollary}
\begin{proof}
The hypotheses of Corollary~\ref{cor:nn-characterization} at the countable
algebraically closed subfields of $k$ hold: Lemma~5.5 of \cite{FHX} by
Theorem~\ref{thm:lpn-local-positivity-instance}, fibre constancy by
Theorem~\ref{thm:gl-fibre-constancy}, and the characteristic-zero envelope
property by Theorem~\ref{thm:kn-smooth-envelope-char-zero}.
\end{proof}
